\section{Definición de la Tecnología HTTPS}

El Protocolo Seguro de Transferencia de Hipertexto (HTTPS, por sus siglas en inglés \textit{Hypertext Transfer Protocol Secure}) constituye la principal tecnología criptográfica empleada para garantizar comunicaciones seguras en la web moderna. Su función esencial radica en proporcionar un canal cifrado, autenticado e íntegro entre un cliente —comúnmente un navegador— y un servidor web. HTTPS se implementa mediante la superposición del protocolo HTTP sobre una capa criptográfica denominada \textit{Transport Layer Security} (TLS), cuyo diseño moderno y estandarización se determinan principalmente en la especificación RFC 8446 \cite{rfc8446}.

Desde una perspectiva funcional, HTTPS permite mitigar amenazas como espionaje, manipulación del tráfico y suplantación de identidad, problemas ampliamente documentados en los análisis formales y empíricos del protocolo TLS \cite{dowling2015, bhargavan2014}. La adopción generalizada de HTTPS ha convertido este estándar en un componente crítico para la infraestructura global de comunicaciones, soportando desde servicios gubernamentales hasta sistemas bancarios, plataformas comerciales y entornos académicos.

En las siguientes subsecciones se expone una definición estructurada, abarcando su origen, sus fundamentos criptográficos y su relevancia en el ecosistema digital contemporáneo.

\subsection{Origen y Evolución de HTTPS}

HTTPS surge como respuesta a las limitaciones de seguridad presentes en el protocolo HTTP original, el cual carecía de mecanismos para evitar la intercepción o la manipulación de datos en tránsito. La introducción de SSL (Secure Sockets Layer) en la década de 1990 constituyó la primera aproximación para resolver estas vulnerabilidades, aunque versiones posteriores demostraron deficiencias críticas. Esto motivó la transición hacia TLS, una familia de protocolos diseñada con mayores garantías formales y sometida a exhaustivas validaciones criptográficas, tales como las realizadas en TLS 1.2 y TLS 1.3 \cite{krawczyk2016, bhargavan2014}.

La última versión ampliamente adoptada, TLS 1.3, representa un avance significativo al reducir la complejidad del \textit{handshake}, mejorar la eficiencia del establecimiento de sesión y reforzar la resistencia frente a ataques contemporáneos, incluyendo técnicas avanzadas de interceptación de tráfico \cite{dowling2015}. La evolución de HTTPS ha sido impulsada no solo por mejoras criptográficas, sino también por la necesidad de garantizar un ecosistema de certificados más confiable, como lo evidencian estudios a gran escala sobre revocación y validez instalada en la infraestructura global de claves públicas (PKI) \cite{liu2015, helme2024}.

\subsection{Concepto Operacional de HTTPS}

En términos operativos, HTTPS no constituye un protocolo independiente, sino un mecanismo de transporte seguro para HTTP. El proceso ocurre mediante el encapsulamiento de las solicitudes y respuestas dentro de un túnel cifrado administrado por TLS. Este túnel se establece mediante un \textit{handshake} criptográfico que cumple tres objetivos fundamentales:

\begin{enumerate}
    \item Autenticación del servidor (y opcionalmente del cliente) mediante certificados digitales X.509 conforme al estándar RFC 5280 \cite{rfc5280}.
    \item Negociación de algoritmos criptográficos robustos para cifrado, integridad y derivación de claves, siguiendo las recomendaciones de uso seguro definidas en RFC 7525 \cite{rfc7525}.
    \item Creación de claves de sesión efímeras empleando generalmente esquemas de intercambio basados en Diffie–Hellman autenticado.
\end{enumerate}

Una vez completado el proceso de negociación, todas las comunicaciones HTTP quedan protegidas mediante cifrado simétrico, lo que garantiza confidencialidad, integridad y autenticidad extremo a extremo.

\subsection{Importancia de HTTPS en la Seguridad de la Información}

HTTPS se ha convertido en un componente esencial para garantizar la seguridad en entornos digitales que dependen de servicios distribuidos, aplicaciones web y arquitecturas basadas en microservicios. En este último caso, la autenticación mutua mediante TLS es un elemento clave para prevenir ataques de manipulación interna y fallas entre servicios, lo cual ha sido documentado en estudios sobre amenazas inter-servicio \cite{sasse2023}. Asimismo, investigaciones sobre la seguridad de inclusiones externas y recursos de terceros han demostrado que HTTPS contribuye a reducir la superficie de ataque asociada a contenido remoto potencialmente malicioso \cite{nikiforakis2012}.

La combinación de estas garantías convierte a HTTPS en una tecnología indispensable para la protección de datos personales, credenciales, sesiones de usuario, transacciones electrónicas y comunicaciones institucionales. Su adopción ya se considera un requisito mínimo en prácticas de seguridad de la información, auditoría digital y cumplimiento normativo.

\subsection{Síntesis de la Definición}

En resumen, HTTPS es un mecanismo criptográfico integral que proporciona seguridad robusta para la comunicación web mediante el uso del protocolo TLS para cifrar, autenticar y asegurar la integridad de los datos. Su diseño ha sido sometido a extensas pruebas formales, su funcionamiento está regulado por estándares internacionales y su adopción constituye un pilar esencial para la protección de la información en el entorno digital contemporáneo.
